\chapter{Stand der Technik}
\label{chap:stand-der-technik}

Auf dem aktuellen Markt der professionellen Veranstaltungstechnik existieren 
bereits verschiedene kommerzielle Lösungen für kabellose und kabelgebundene LED-Röhrensysteme. 
Nachfolgend werden ausgewählte, marktführende Produkte 
hinsichtlich ihres Funktionsumfangs, ihres Einsatzbereiches und ihrer 
Kostenstrukturen analysiert, um den aktuellen Stand der Technik abzubilden 
und Lücken im bestehenden Angebot aufzuzeigen. 

\section{Kommerzielle LED-Leuchtröhren}
\label{sec:kommerzielle-led-leuchtroehren}

\subsection{Premium-Segment: Astera Titan Tube}
Als gegenwärtiger Marktführer im Bereich autarker und per Funk steuerbarer 
LED-Röhren gilt das Unternehmen Astera mit seinem Modell \textit{Titan Tube}. 
Dabei handelt es sich um eine einen Meter lange LED-Röhre, die auf einer 
RGBWMA-Pixelarchitektur (Red, Green, Blue, Mint, Amber) basiert. Durch die 
Hinzunahme der Mint- und Amber-Dioden ermöglicht das System eine
Anpassung des Weißabgleichs. 
Aufgrund dieser Eigenschaften eignet sich die Röhre neben dem Einsatz im 
Eventbereich auch für professionelle Fotografie- und Videografieproduktionen.

Für die kabellose Kommunikation unterstützt das System Wireless DMX
im 2,4-GHz-Frequenzband. Für den mobilen 
Einsatz verfügt die Röhre über einen integrierten Akku, der bei maximaler 
Helligkeit und Ausgabe von reinem Weiß eine spezifizierte Laufzeit von etwa 
1 Stunde und 45 Minuten erreicht. Die Schutzklasse IP65 zertifiziert das 
Gehäuse zusätzlich als staubdicht und strahlwassergeschützt, was eine 
zuverlässige Nutzung im Außenbereich garantiert.

Darüber hinaus unterstützt die Titan Tube eine
kabelgebundene Pixelsteuerung über die etablierten Protokolle DMX512 sowie 
Art-Net. Dieser Betriebsmodus ist aufgrund der hohen Datenrate jedoch an den Betrieb
über Kabel gebunden. Die hohe Funktionsdichte 
und Verarbeitungsqualität spiegelt sich in der Preisstruktur 
wider: Eine einzelne Röhre kostet ca. 800€, während ein Komplettset 
aus acht Röhren inklusive Transportbox und Ladegerät mit rund 
7.000€ kalkuliert ist.

\subsection{Kostengünstige Alternativen: BeamZ Zelos}
Als preisgünstigere Alternative im professionellen Segment lassen sich 
Systeme wie das Modell \textit{Zelos} des Herstellers BeamZ identifizieren. 
Diese Systeme richten sich mit einem Gesamtpreis von etwa 3.000€ für 
ein 8er-Set (was rund 400€ pro Röhre entspricht) primär an Anwendungsfälle 
im Veranstaltungsbetrieb. 

Zwar verfügt auch dieses System über eine drahtlose DMX-Ansteuerung im 
2,4-GHz-Band, jedoch bindet es den Einsatzbereich an einen reinen Netzbetrieb.
Ein weitaus signifikanterer Nachteil gegenüber dem 
Premiumsegment ist das Fehlen von Datenschnittstellen für eine kabelgebundene 
Pixelsteuerung. Die Ansteuerung beschränkt sich somit auf die grundlegende 
Funkbedienung ohne die Systemflexibilität für ein hochauflösendes, 
drahtgebundenes Pixel-Mapping (etwa via Art-Net).

\section{Wissenschaftliche Forschung und verwandte Konzepte}
\label{sec:wissenschaftliche-forschung}

Neben den am Markt verfügbaren Produkten existieren wissenschaftliche Arbeiten,
die sich mit der drahtlosen Übertragung von Lichtsteuerdaten befassen. Der
folgende Abschnitt ordnet diese Ansätze ein und grenzt sie gegenüber dem in
dieser Arbeit verfolgten Konzept ab.

\subsection{Drahtlose Kommunikationsprotokolle in der Beleuchtungstechnik}
\label{subsec:forschung-kommunikation}

Die Ablösung der kabelgebundenen DMX512-Übertragung durch eine Funkstrecke ist
bereits Gegenstand wissenschaftlicher Untersuchungen. Yang et al. stellen mit
\textit{WiDMX512} ein System vor, das die Nachteile der klassischen
Bus-Topologie explizit zum Ausgangspunkt nimmt: Die kabelgebundene Verkettung
von Master und Slaves erhöhe die Systemkomplexität insbesondere bei großen
Leitungslängen, schränke Mobilität und Skalierbarkeit ein und verursache
zusätzlichen Material- und Personalaufwand \cite{yang2015real}. Diese
Problemanalyse deckt sich unmittelbar mit der in
Abschnitt~\ref{sec:problemstellung} dargelegten Ausgangslage. Sie belegt, dass
es sich nicht um ein Einzelfallproblem der Veranstaltungspraxis handelt, sondern
um eine Eigenschaft des Protokollstandards.

Technisch setzt WiDMX512 auf Transceiver-Module des Typs Texas Instruments
CC1110 im 915-MHz-ISM-Band bei einer Übertragungsrate von 250\,kbit/s. Da die
Übertragung eines vollständigen DMX512-Pakets von 512\,Byte bei dieser Datenrate
31{,}2\,ms beansprucht und damit die im Standard vorgesehene Aktualisierungszeit
von 22{,}7\,ms überschreitet, teilen die Autoren jedes Frame in zwei Hälften zu
je 256\,Byte und übertragen diese über zwei parallel betriebene Transceiver auf
getrennten Kanälen. Empfängerseitig wird das Frame über ein Doppelpufferverfahren
wieder zusammengesetzt. Oszilloskopmessungen weisen eine Verarbeitungszeit von
15{,}2\,ms auf der Senderseite und von 22{,}7\,ms bis zum vollständigen Eintreffen
am Leuchtmittel nach, womit das System die Zeitvorgaben des Standards einhält
\cite{yang2015real}.

Ein zweiter Beitrag des Systems liegt in seiner bidirektionalen Auslegung.
Während DMX512 unidirektional arbeitet und der Sender folglich keine Rückmeldung
darüber erhält, ob ein Gerät korrekt adressiert wurde, quittieren die Slaves in
WiDMX512 jedes empfangene Paket. Kanal- und Adresszuweisung erfolgen über ein
Handshake-Verfahren mit bis zu drei Wiederholungen je Anfrage. Für Geräte, die
außerhalb der Funkreichweite liegen, durch Hindernisse abgeschattet werden oder
auf einem abweichenden Kanal betrieben werden, schlagen die Autoren eine gemischte
Topologie vor: Nicht erreichbare Slaves werden über ein Kabel an einen erreichbaren
Slave angebunden und von diesem als Relais bedient \cite{yang2015real}.

Einen abweichenden Transportweg untersucht Newton, der die Übertragung des
Art-Net-Protokolls über handelsübliche WLAN-Router im Kontext der
Beleuchtungssteuerung betrachtet \cite{1554184}. Dieser Ansatz nutzt vorhandene
Netzwerkinfrastruktur, bindet die Übertragung damit jedoch an den vollständigen
IEEE-802.11-Stack einschließlich Assoziierung und Access-Point-Verwaltung.

Aus beiden Arbeiten ergeben sich für die vorliegende Arbeit vier Abgrenzungen,
die weniger als Qualitätsunterschied denn als Folge abweichender Randbedingungen
zu verstehen sind.

\begin{description}
  \item[Frequenzbereich] Das von Yang et al. genutzte 915-MHz-Band steht in
    Europa nicht in gleicher Weise anmeldefreie zur Verfügung wie
    in Amerika oder in Südkorea, hier ist der Bereich von 863 bis
    870\,MHz vorgesehen. Der Ansatz ist damit nicht unmittelbar auf den
    angestrebten Einsatzort übertragbar. Die vorliegende Arbeit nutzt stattdessen
    das weltweit verfügbare 2,4-GHz-Band
    (Abschnitt~\ref{subsec:drahtlose-uebertragung}).
  \item[Datenrate und Redundanz] WiDMX512 muss zwei Transceiver parallel
    betreiben, um ein einziges DMX-Frame fristgerecht zu übertragen, und verfügt
    damit über keinerlei Übertragungsreserve. Die deutlich höhere Bruttodatenrate
    des hier eingesetzten Übertragungsverfahrens erlaubt es hingegen, jeden
    Datenstand mehrfach zu übertragen. Die in Kapitel~\ref{chap:evaluierung}
    untersuchte Redundanzstrategie ist folglich keine Notlösung, sondern eine
    erst durch die Wahl der Übertragungsschicht ermöglichte Entwurfsoption.
  \item[Quittierung gegenüber Broadcast] Die Bestätigung jedes Pakets verschafft
    WiDMX512 Zustellsicherheit und eine Rückmeldung über den Konfigurationszustand
    der Empfänger, erzeugt jedoch einen Aufwand, der mit der Zahl der Empfänger
    wächst. Die vorliegende Arbeit verzichtet zugunsten einer konstanten
    One-to-Many-Latenz auf den Rückkanal und kompensiert die fehlende Bestätigung
    durch mehrfache Aussendung
    (Abschnitt~\ref{subsec:drahtlose-uebertragung}). Die Adressierung erfolgt
    nicht über ein Funkprotokoll, sondern lokal am Gerät
    (Abschnitt~\ref{subsec:benutzerfreundlichkeit}).
\end{description}

Ergänzend ist anzumerken, dass die Leistungsbewertung in \cite{yang2015real} auf
den Nachweis der Timing-Konformität beschränkt bleibt. Angaben zu Reichweite,
Paketverlust oder zum Verhalten unter Fremdfunklast werden nicht erhoben, obwohl
gerade diese Größen für den Veranstaltungsbetrieb maßgeblich sind. Die in
Kapitel~\ref{chap:evaluierung} durchgeführten Messreihen setzen an dieser Stelle
an.

\subsection{Synchronisationskonzepte in verteilten Echtzeitsystemen}
\label{subsec:forschung-synchronisation}

Werden Lichteffekte auf mehrere räumlich verteilte Geräte aufgeteilt, muss deren
Ausgabe für den Betrachter gleichzeitig erfolgen. Der klassische Lösungsweg besteht
darin, die lokalen Uhren der Knoten über ein Synchronisationsprotokoll anzugleichen.
Elson et al.\ verfolgen mit der \textit{Reference-Broadcast Synchronization} einen
Ansatz, der für die vorliegende Arbeit unmittelbar aufschlussreich ist: Statt einen
Zeitstempel vom Sender zum Empfänger zu übertragen, nutzen sie den Umstand, dass ein
auf der Bitübertragungsschicht ausgesendeter Broadcast von allen Empfängern nahezu
gleichzeitig registriert wird. Die nichtdeterministischen Verzögerungen auf der
Senderseite, insbesondere die Wartezeit auf den Kanalzugriff wirken auf alle
Empfänger gleichermaßen und fallen aus der Betrachtung heraus. Mit handelsüblicher
802.11-Hardware erreichen die Autoren auf diese Weise eine Übereinstimmung der Uhren
von 1,85\,µs bei einer Standardabweichung von 1,28\,µs \cite{elson2002}.

Das entwickelte System verzichtet auf ein Synchronisationsprotokoll und nutzt
stattdessen dieselbe Eigenschaft des Mediums unmittelbar: Alle Röhren empfangen
denselben Broadcast innerhalb desselben Übertragungsfensters
(Abschnitt~\ref{subsec:synchronisation-roehem}). Zulässig ist dieser Verzicht, weil die
hier geforderte Genauigkeit nicht im Mikrosekundenbereich liegt, sondern aus der
Bewegungswahrnehmungsschwelle folgt und mit rund 42\,ms
(Abschnitt~\ref{subsec:echtzeitverhalten}) vier Größenordnungen darüber liegt. Die
Arbeit von Elson et al.\ zeigt damit, dass es sich nicht um eine bloße Vereinfachung
handelt, sondern um die Nutzung einer dokumentierten Eigenschaft von Broadcast-Medien.
Nicht gewonnen wird dadurch allerdings eine gemeinsame Zeitbasis zwischen Sender und
Empfänger, deren Fehlen ist der Grund, weshalb die Ende-zu-Ende-Latenz in dieser Arbeit
nicht messbar ist (Abschnitt~\ref{sec:grenzen}).

\subsection{Echtzeit-Signalverarbeitung auf Microcontrollern}
\label{subsec:forschung-signalverarbeitung}

Eine dritte Gruppe von Arbeiten betrifft nicht die Übertragung selbst, sondern die
Frage, ob der empfangende Mikrocontroller die Daten fristgerecht verarbeiten kann.
Arm et al.\ untersuchen hierzu das Laufzeitverhalten von FreeRTOS auf einem
ESP32-WROOM-32 mit Xtensa-Dual-Core-Prozessor und messen unter anderem die Dauer von
Kontextwechseln sowie den Jitter der Task-Perioden. Ihr zentrales Ergebnis ist, dass
Kontextwechselzeit und Jitter deutlich ansteigen, sobald die Ausführung eines Tasks auf
den jeweils anderen Kern verlagert wird \cite{arm2022}.

Für die vorliegende Arbeit folgen daraus zwei Konsequenzen. Erstens ergibt sich die
Entwurfsvorgabe, zeitkritische Abläufe fest einem Kern zuzuordnen, statt sie zwischen
den Kernen wandern zu lassen, die in Abschnitt~\ref{subsec:nebenlaeufigkeit}
beschriebene Aufteilung: Funkstack auf dem Protokollkern, Effektberechnung und
LED-Ausgabe auf dem Applikationskern, Übergabe über einen überschreibenden Puffer
setzt genau dies um. Zweitens zeigt die Arbeit, dass Jitter nicht ausschließlich der
Funkstrecke zuzurechnen ist, sondern auch auf dem Mikrocontroller selbst entsteht. Der
in Kapitel~\ref{chap:evaluierung} erfasste Jitter enthält diesen Anteil folglich mit
und ist als Eigenschaft der gesamten Verarbeitungskette zu lesen.

\section{Identifizierte Lücken im Stand der Technik}
\label{sec:luecken-stand-der-technik}

Die Marktanalyse zeigt eine strukturelle Lücke zwischen funktional umfassenden, 
jedoch kostenintensiven Premiumsystemen und preisgünstigen, zugleich deutlich 
reduzierten Alternativen. Während Systeme wie die Astera \textit{Titan Tube} 
drahtlose Betriebsfähigkeit und kabelgebundene Pixelsteuerung in hoher Integrationsqualität 
vereinen, sind vergleichbare Funktionsumfänge im unteren Preissegment nicht verfügbar.

Die drahtlose Übertragung von DMX-Steuerdaten ist demgegenüber als Einzelproblem
wissenschaftlich bereits behandelt worden
(Abschnitt~\ref{subsec:forschung-kommunikation}). Die Forschungslücke besteht
folglich nicht darin, dass eine funkgestützte Ansteuerung grundsätzlich
unerforscht wäre.

Sie liegt vielmehr im Grad der Integration. Die betrachteten Arbeiten ersetzen
jeweils nur die Übertragungsstrecke zwischen einem Steuergerät und bestehenden
Beleuchtungsgeräten, während das Leuchtmittel selbst unverändert bleibt und
weiterhin kabelgebunden versorgt und angesteuert wird. Ein System, das
Funkempfang und Energieversorgung innerhalb des
Leuchtmittels vereint und dadurch die Signalverkabelung vollständig entfallen
lässt, wird in den betrachteten Arbeiten nicht beschrieben. Der Ansatz von Yang
et al. verdeutlicht diese Grenze besonders, da dort Reichweiten- und
Abschattungsprobleme durch zusätzliche Kabelverbindungen zwischen den Empfängern
gelöst werden \cite{yang2015real}.

Hinzu kommt die Frage der Verfügbarkeit: Die kommerziellen Systeme sind
proprietär, die wissenschaftlichen Arbeiten beschreiben Prototypen ohne
veröffentlichte Umsetzung. Eine reproduzierbar und transparent dokumentierte
Open-Source-Lösung, die zugleich an etablierte Steuerungsumgebungen
anschlussfähig bleibt, ist im betrachteten Anwendungsfeld bislang nicht
etabliert. Die Neuheit der vorliegenden Arbeit besteht folglich in der
Entwicklung einer offenen Funkarchitektur, die Übertragung und Leuchtmittel als
Gesamtsystem begreift.

